\section{Chuyển đổi từ Mô hình Quan niệm sang Lược đồ Logic}

Dựa trên sơ đồ thực thể kết hợp (ERD) và các quy tắc nghiệp vụ đã phân tích ở Chương 2, tiến hành chuyển đổi mô hình thực thể quan niệm sang mô hình quan hệ (Relational Schema). Quá trình này được thực hiện thông qua việc áp dụng nghiêm ngặt các thuật toán ánh xạ ER sang quan hệ.

\subsection{Các quy tắc chuyển đổi áp dụng}
Quá trình ánh xạ từ ERD sang Lược đồ Logic tuân thủ các quy tắc cốt lõi sau:

\begin{itemize}
    \item \textbf{Quy tắc 1 (Ánh xạ Thực thể mạnh):} Mỗi thực thể (Entity) mạnh trong ERD được chuyển thành một quan hệ (Bảng - Table) trong Lược đồ Logic. Các thuộc tính đơn trở thành các cột (Columns) của bảng. Khóa định danh của thực thể trở thành Khóa chính (Primary Key - PK).
    
    \item \textbf{Quy tắc 2 (Ánh xạ mối kết hợp 1:N):} Đối với các mối kết hợp một-nhiều (ví dụ: \texttt{Merchant} - \texttt{Campaign}, \texttt{Customer} - \texttt{Credit}), đưa khóa chính của bảng phía "1" sang bảng phía "N" làm khóa ngoại (Foreign Key - FK).
    
    \item \textbf{Quy tắc 3 (Ánh xạ mối kết hợp M:N):} Tạo một quan hệ trung gian cho mối kết hợp nhiều-nhiều (ví dụ Customer và Merchant thông qua Credit). Đưa khóa của hai thực thể vào làm khóa ngoại. Trong CREDIT, chọn CreditID làm PK và đặt UNIQUE(CustomerID, MerchantID) cho khóa nghiệp vụ; không ghép CreditID vào cặp khóa này.
    \item \textbf{Quan hệ tự tham chiếu:} Đưa ManagerID nullable vào USER làm FK tham chiếu UserID, giữ nguyên quan hệ quản lý của ERD.
\end{itemize}

\subsection{Danh sách Lược đồ Quan hệ sơ bộ}
Dựa trên ERD của OctaPoint CaaS, áp dụng các quy tắc trên, ta thu được tập hợp các lược đồ quan hệ sơ bộ như sau. (Quy ước: Thuộc tính gạch chân là Khóa chính - PK, thuộc tính in nghiêng là Khóa ngoại - FK).

\begin{description}
    \item[\texttt{Tenant}] (\underline{TenantID}, TenantName, SystemConfig)
    \item[\texttt{Merchant}] (\underline{MerchantID}, \textit{TenantID}, BusinessName, Status)
    \item[\texttt{Role}] (\underline{RoleID}, RoleName)

    \item[\texttt{User}] (\underline{UserID}, \textit{MerchantID}, Email, PasswordHash, \textit{RoleID}, \textit{ManagerID})
    \item[\texttt{Campaign}] (\underline{CampaignID}, \textit{MerchantID}, CampaignName, StartDate, EndDate, Status)
    \item[\texttt{RewardRule}] (\underline{RuleID}, \textit{CampaignID}, ActionType, PointsMultiplier)
    \item[\texttt{User}] (\underline{UserID}, \textit{MerchantID}, Email, PasswordHash, \textit{RoleID}, \textit{ManagerID})    
    \item[\texttt{Campaign}] (\underline{CampaignID}, \textit{MerchantID}, CampaignName, StartDate, EndDate, Status)
    \item[\texttt{RewardRule}] (\underline{RuleID}, \textit{CampaignID}, ActionType, PointsMultiplier) 
    \item[\texttt{Customer}] (\underline{CustomerID}, WalletAddress, Email)
    \item[\texttt{Credit}] (\underline{CreditID}, \textit{CustomerID}, \textit{MerchantID}, AvailableBalance)
    \item[\texttt{CreditTransaction}] (\underline{TransactionID}, \textit{CreditID}, \textit{CampaignID}, Amount, TransactionType, SuiTxDigest, CreatedAt)
    \item[\texttt{APICredential}] (\underline{KeyID}, \textit{MerchantID}, ApiKey, SecretHash)
\end{description}
